\section{Task Scheduler}
\label{sec:heap-task-scheduler}

\problemstatement{We are given an array $\textit{tasks}$, where each
element is a task type (e.g.\ a character), and an integer $n$, the
\textbf{cooldown}: the same task type cannot be executed again until at
least $n$ units of time have passed since its last execution (the CPU may
sit \textbf{idle} for a unit of time if no eligible task is available).
Find the \textbf{minimum total time} needed to execute all tasks.}

\begin{patternbox}
\emph{``Always schedule whichever task type currently has the most
remaining occurrences (subject to a cooldown), repeatedly''} is a
perpetually-resorted greedy priority -- exactly the
\emph{``schedule/group with a cost or cooldown that always prefers the
largest remaining count first''} keyword pattern from
Section~\ref{sec:heap-intro}. Scheduling the most frequent task type first,
whenever it is eligible, spreads out the most demanding task type as early
and as evenly as possible, minimizing the idle time it would otherwise
force later.
\end{patternbox}

\subsection*{Understanding the problem carefully}

Intuitively, the task type with the highest remaining count is the most
constraining: if we delay it, we are more likely to be forced into idle
slots later, because its cooldown keeps recurring. So at every time step,
we should greedily execute whichever \emph{eligible} task type currently
has the largest remaining count. A max-heap keyed by remaining count gives
us this ``most frequent remaining task'' in $O(\log d)$ time, where $d$ is
the number of distinct task types (at most $26$ for lowercase letters, so
effectively a small constant).

\begin{ideabox}[Greedy Choice, Powered by a Max-Heap]
Count the frequency of each task type and push these counts into a
max-heap. Process time in chunks of $n+1$ consecutive units (one full
cooldown cycle): within each chunk, repeatedly pop the most frequent
remaining task type, decrement its count, and execute it in the next time
slot; if its count is still positive after decrementing, set it aside
(it cannot be reused within this same cooldown cycle) to be pushed back
into the heap only after the chunk completes. If both the heap and the
set-aside list are empty before filling all $n+1$ slots of the current
chunk, stop immediately (no trailing idle time is needed); otherwise, any
unfilled slots in the chunk count as idle time.
\end{ideabox}

\subsection*{Why processing in cooldown-sized chunks is correct}

Within a single chunk of $n+1$ consecutive time slots, no task type can
appear twice (that is exactly what the cooldown $n$ forbids -- a task
executed at the start of a chunk cannot be executed again until at least
$n$ slots later, which is precisely the boundary of the next chunk). So a
chunk can hold at most one execution of each distinct task type, and
greedily filling each chunk with the currently most-frequent remaining
task types (up to $\min(d, n+1)$ of them, where $d$ is how many distinct
types still have remaining work) never wastes an opportunity: any
slot left empty in a chunk after all distinct remaining types have been
used \emph{must} be idle, because no task is eligible. Choosing the
\emph{most frequent} remaining types to fill each chunk (rather than any
other selection) is what minimizes the number of future chunks that will
be forced to leave slots idle, since it depletes the most constraining
task types fastest, evening out the workload as quickly as possible.

\subsection*{Algorithm}

\begin{enumerate}
  \item Count the frequency of each task type; push all counts into a
        max-heap.
  \item Initialize $\textit{time} \gets 0$.
  \item While the heap is not empty:
  \begin{enumerate}
    \item $\textit{tempList} \gets$ empty list (tasks used this cycle,
          temporarily held out of the heap).
    \item For $i \gets 0$ to $n$ ($n+1$ slots in this cycle):
    \begin{enumerate}
      \item If the heap is not empty: pop the largest count, decrement
            it by $1$; if still $>0$, append it to $\textit{tempList}$.
      \item $\textit{time} \gets \textit{time}+1$.
      \item If the heap is empty and $\textit{tempList}$ is empty: break
            out of this inner loop early (no trailing idle needed).
    \end{enumerate}
    \item Push every count in $\textit{tempList}$ back onto the heap.
  \end{enumerate}
  \item Return \textit{time}.
\end{enumerate}

\begin{algorithm}[H]
\caption{Task Scheduler}
\begin{algorithmic}[1]
\Procedure{LeastInterval}{\textit{tasks}, $n$}
  \State count frequencies; push all into max-heap \textit{heap}
  \State $\textit{time} \gets 0$
  \While{\textit{heap} is not empty}
    \State $\textit{temp} \gets$ empty list
    \For{$i \gets 0$ to $n$}
      \If{\textit{heap} is not empty}
        \State $c \gets$ pop maximum; $c \gets c - 1$
        \If{$c > 0$} append $c$ to \textit{temp} \EndIf
      \EndIf
      \State $\textit{time} \gets \textit{time} + 1$
      \If{\textit{heap} is empty \textbf{and} \textit{temp} is empty}
        \State \textbf{break}
      \EndIf
    \EndFor
    \State push every count in \textit{temp} back onto \textit{heap}
  \EndWhile
  \State \Return \textit{time}
\EndProcedure
\end{algorithmic}
\end{algorithm}

\subsection*{Dry run}

\dryrun{Take $\textit{tasks} = [A,A,A,B,B,B]$, $n = 2$.}

Counts: $A{:}3, B{:}3$. Heap $=\{3_A, 3_B\}$.

\begin{itemize}
  \item \textbf{Cycle 1} ($3$ slots, since $n+1=3$): pop $A$ (or $B$;
        ties broken arbitrarily), decrement to $2$, set aside. Slot
        used, $\textit{time}=1$. Pop $B$, decrement to $2$, set aside.
        Slot used, $\textit{time}=2$. Heap now empty, temp has
        $\{2_A, 2_B\}$ but temp is not empty, so we still fill the third
        slot: heap is empty, so this slot is idle. $\textit{time}=3$.
        Push $2_A, 2_B$ back: heap $=\{2_A, 2_B\}$.
  \item \textbf{Cycle 2}: pop $A$ (decrement to $1$, set aside),
        $\textit{time}=4$. Pop $B$ (decrement to $1$, set aside),
        $\textit{time}=5$. Third slot: heap empty, temp not empty, idle.
        $\textit{time}=6$. Push $1_A, 1_B$ back.
  \item \textbf{Cycle 3}: pop $A$ (decrement to $0$, do not set aside),
        $\textit{time}=7$. Pop $B$ (decrement to $0$, do not set aside),
        $\textit{time}=8$. Now heap is empty and temp is empty: break
        early, no third idle slot needed.
\end{itemize}

Total time is $8$, matching the well-known schedule
$A,B,\text{idle},A,B,\text{idle},A,B$.

\subsection*{C++ implementation}

\begin{lstlisting}
#include <bits/stdc++.h>
using namespace std;

int leastInterval(vector<char>& tasks, int n) {
    unordered_map<char,int> freq;
    for (char t : tasks) freq[t]++;

    priority_queue<int> maxHeap;
    for (auto& [ch, c] : freq) maxHeap.push(c);

    int time = 0;
    while (!maxHeap.empty()) {
        vector<int> temp;
        int slots = n + 1;

        for (int i = 0; i < slots; i++) {
            if (!maxHeap.empty()) {
                int c = maxHeap.top() - 1;
                maxHeap.pop();
                if (c > 0) temp.push_back(c);
            }
            time++;
            if (maxHeap.empty() && temp.empty()) break;
        }
        for (int c : temp) maxHeap.push(c);
    }
    return time;
}

int main() {
    vector<char> tasks = {'A','A','A','B','B','B'};
    cout << "Minimum time: " << leastInterval(tasks, 2) << endl;
    return 0;
}
\end{lstlisting}

\begin{complexitybox}
\textbf{Time Complexity.} There are at most $d \le 26$ distinct task
types, so the heap has size $O(1)$ from a complexity standpoint, and the
total number of pop/push operations is bounded by the total task count
$N$ (each task is popped and re-pushed $O(1)$ times across all cycles in
the typical analysis), giving
\[
  O(N \log d) = O(N),
\]
effectively linear in the number of tasks.
\textbf{Space Complexity.} $O(d)$ for the frequency map and heap.
\end{complexitybox}

\begin{takeawaybox}
A closed-form formula also exists for this specific problem
($\textit{answer} = \max(N, (\textit{maxFreq}-1)(n+1) + \textit{numTasksWithMaxFreq})$),
but the heap-based simulation shown here generalizes better: it directly
extends to variants with per-task-type costs, partial cooldowns, or other
constraints where a clean closed form may not exist, while the
underlying ``always schedule the currently most-constraining option'' greedy
idea remains unchanged.
\end{takeawaybox}
